\chapter{Files, Paths, and the Filesystem}\label{ch:filesystem}

\index{filesystem}
\index{file I/O}

Every program eventually needs to talk to the outside world, and for most
programs, that world is the filesystem. Whether you are reading a configuration
file at startup, writing a log, or processing a batch of images, Lattice gives
you a set of focused, no-surprises functions that map cleanly onto the
underlying operating system. There are no deep class hierarchies to learn and
no ``stream'' abstractions to wire together---just functions that take paths
and return values.

In this chapter we will start with the three workhorses of text file I/O:
\lstinline|read_file|, \lstinline|write_file|, and \lstinline|append_file|.
From there we will explore the broader \texttt{fs} family of functions for
querying and manipulating the filesystem, the \texttt{path} helpers for
constructing and dissecting file paths, temporary files and directories, and
finally the \lstinline|Buffer| type for working with raw binary data.

% ---------------------------------------------------------------------------
\section{read\_file, write\_file, append\_file}\label{sec:read-write-append}
% ---------------------------------------------------------------------------

\index{read\_file}
\index{write\_file}
\index{append\_file}

Let's start with the most common operation in scripting: reading a file into a
string.

\begin{lstlisting}[language=Lattice, caption={Reading and printing a file}]
let contents = read_file("config.toml")
print(contents)
\end{lstlisting}

\lstinline|read_file| takes a path as a \lstinline|String| and returns the
entire file contents as a \lstinline|String|. If the file does not exist or
cannot be opened, it returns \lstinline|nil|. Under the hood, the
implementation in \filename{src/builtins.c} opens the file with
\texttt{fopen}, seeks to the end to determine the length, allocates exactly
enough memory, and reads the contents in one pass.

\begin{notebox}{Encoding Assumption}
\lstinline|read_file| assumes the file contains valid UTF-8 text. If you need
to read arbitrary binary data---an image, a compiled bytecode file, a
protocol buffer---use \lstinline|read_file_bytes| instead, which returns a
\lstinline|Buffer|. We cover Buffers in \cref{sec:buffers}.
\end{notebox}

Writing is the mirror operation:

\begin{lstlisting}[language=Lattice, caption={Writing a string to a file}]
let report = "Scan completed: 0 issues found.\n"
write_file("report.txt", report)
\end{lstlisting}

\lstinline|write_file| creates the file if it does not exist, or
\emph{truncates} it if it does. It returns \lstinline|true| on success and
\lstinline|false| on failure. This is a full overwrite---if you need to add
lines to an existing file, use \lstinline|append_file|:

\begin{lstlisting}[language=Lattice, caption={Appending to a log file}]
fn log_event(message: String) {
    let timestamp = time_format(time_now(), "%Y-%m-%d %H:%M:%S")
    let line = "[" + timestamp + "] " + message + "\n"
    append_file("app.log", line)
}

log_event("Server started on port 8080")
log_event("Accepted connection from 192.168.1.42")
\end{lstlisting}

\lstinline|append_file| opens the file in append mode (\texttt{"a"} in C
terms), so writes are guaranteed to go to the end of the file even if another
process is writing concurrently. Like \lstinline|write_file|, it creates the
file if it does not exist.

\begin{warnbox}{write\_file Truncates!}
A common mistake is to use \lstinline|write_file| when you meant
\lstinline|append_file|. If your program writes a log file using
\lstinline|write_file| in a loop, each iteration will obliterate everything
the previous iteration wrote. When in doubt, ask yourself: ``Do I want to
replace or extend?''
\end{warnbox}

\subsection{A Complete Read-Modify-Write Cycle}\label{sec:read-modify-write}

A typical pattern combines all three operations. Suppose we have a simple
counter file that stores a single integer:

\begin{lstlisting}[language=Lattice, caption={Incrementing a persistent counter}]
fn increment_counter(path: String) -> Int {
    flux count = 0

    if file_exists(path) {
        let text = read_file(path)
        count = parse_int(text.trim())
    }

    count += 1
    write_file(path, to_string(count))
    return count
}

let visits = increment_counter("/tmp/visit_count.txt")
print("Visit #" + to_string(visits))
\end{lstlisting}

Notice how we check \lstinline|file_exists| before reading. This is a
deliberate choice---\lstinline|read_file| returns \lstinline|nil| on a missing
file, but checking explicitly makes the intent clearer and lets us set a
sensible default.

\begin{tipbox}{Prefer try/catch for Robust Code}
For production code, wrap file operations in \lstinline|try|/\lstinline|catch|
blocks (see \cref{ch:error-handling}). The examples in this chapter omit
error handling for clarity, but real programs should always account for
missing files, permission errors, and full disks.
\end{tipbox}

% ---------------------------------------------------------------------------
\section{The fs Module}\label{sec:fs-module}
% ---------------------------------------------------------------------------

\index{file\_exists}
\index{list\_dir}
\index{mkdir}
\index{glob}
\index{stat}
\index{is\_dir}
\index{is\_file}
\index{delete\_file}
\index{copy\_file}
\index{rename}
\index{rmdir}
\index{chmod}
\index{realpath}
\index{file\_size}

Beyond reading and writing file contents, Lattice exposes a family of
functions for querying and manipulating the filesystem itself. These are
registered as top-level built-in functions---no import is required.

\subsection{Checking Existence: file\_exists, is\_file, is\_dir}
\label{sec:checking-existence}

\begin{lstlisting}[language=Lattice, caption={Checking existence and type}]
let path = "/etc/hosts"

print(file_exists(path))   // true
print(is_file(path))       // true
print(is_dir(path))        // false
print(is_dir("/etc"))      // true
\end{lstlisting}

\lstinline|file_exists| returns \lstinline|true| if \emph{anything} exists at
the given path---file, directory, symlink, or device node. If you need to
distinguish between files and directories, use \lstinline|is_file| and
\lstinline|is_dir|. All three delegate to the POSIX \texttt{stat} system call
internally (see \filename{src/fs\_ops.c}).

\subsection{Listing Directories: list\_dir}\label{sec:list-dir}

\begin{lstlisting}[language=Lattice, caption={Listing directory entries}]
let entries = list_dir(".")
for entry in entries {
    print(entry)
}
\end{lstlisting}

\lstinline|list_dir| returns an array of strings---one per entry in the
directory, excluding the special \texttt{.} and \texttt{..} entries. The
entries are file \emph{names}, not full paths. If you need full paths, combine
with \lstinline|path_join|:

\begin{lstlisting}[language=Lattice, caption={Building full paths from directory entries}]
let dir = "/var/log"
let entries = list_dir(dir)
for name in entries {
    let full_path = path_join(dir, name)
    let info = stat(full_path)
    print(full_path + " (" + info["type"] + ", " + to_string(info["size"]) + " bytes)")
}
\end{lstlisting}

\subsection{Creating and Removing Directories: mkdir, rmdir}
\label{sec:mkdir-rmdir}

\begin{lstlisting}[language=Lattice, caption={Creating and removing directories}]
mkdir("output")
write_file("output/results.csv", "name,score\nAlice,97\n")

// Later, clean up:
delete_file("output/results.csv")
rmdir("output")
\end{lstlisting}

\lstinline|mkdir| creates a single directory with permissions \texttt{0755}.
It does \emph{not} create intermediate directories---if you need
\texttt{mkdir -p} behavior, you will need to create each level yourself:

\begin{lstlisting}[language=Lattice, caption={Creating nested directories}]
fn mkdir_p(path: String) {
    let parts = path.split("/")
    flux current = ""
    for part in parts {
        if part == "" {
            current = "/"
            continue
        }
        current = path_join(current, part)
        if !file_exists(current) {
            mkdir(current)
        }
    }
}

mkdir_p("output/reports/2024/q4")
\end{lstlisting}

\lstinline|rmdir| removes an \emph{empty} directory. Attempting to remove a
non-empty directory will produce an error. If you need to recursively remove a
directory tree, you must walk it yourself, deleting files first and
directories bottom-up.

\subsection{Globbing: glob}\label{sec:glob}

\index{glob}

When you need to find files matching a pattern, \lstinline|glob| is your
friend:

\begin{lstlisting}[language=Lattice, caption={Finding files with glob patterns}]
// Find all Lattice source files in the current directory
let sources = glob("*.lat")
for path in sources {
    print("Found: " + path)
}

// Find all JPEG images in a photos directory
let photos = glob("photos/*.jpg")
print("Found " + to_string(photos.len()) + " photos")
\end{lstlisting}

The pattern syntax follows the POSIX \texttt{glob(3)} convention:
\texttt{*} matches any sequence of characters within a path component,
\texttt{?} matches a single character, and \texttt{[abc]} matches character
classes. On systems that support it, brace expansion
(\texttt{\{jpg,png,gif\}}) is also available.

\begin{lstlisting}[language=Lattice, caption={Brace expansion in glob patterns}]
// Find all image files (on systems with GLOB_BRACE support)
let images = glob("assets/*.{jpg,png,gif,webp}")
\end{lstlisting}

If no files match, \lstinline|glob| returns an empty array---it does not treat
the absence of matches as an error. You can see this behavior in
\filename{src/fs\_ops.c}, where \texttt{GLOB\_NOMATCH} returns \texttt{NULL}
with a count of zero rather than setting the error string.

\subsection{File Metadata: stat and file\_size}\label{sec:stat}

\index{stat}
\index{file\_size}

The \lstinline|stat| function returns a map with detailed metadata about a
path:

\begin{lstlisting}[language=Lattice, caption={Inspecting file metadata with stat}]
let info = stat("README.md")
print("Type: " + info["type"])           // "file"
print("Size: " + to_string(info["size"]) + " bytes")
print("Modified: " + to_string(info["mtime"]))  // epoch milliseconds
print("Mode: " + to_string(info["mode"]))        // permission bits
\end{lstlisting}

The returned map contains five keys:

\begin{center}
\begin{tabular}{lll}
\toprule
\textbf{Key} & \textbf{Type} & \textbf{Description} \\
\midrule
\lstinline|"type"| & \lstinline|String| & \texttt{"file"}, \texttt{"dir"}, \texttt{"symlink"}, or \texttt{"other"} \\
\lstinline|"size"| & \lstinline|Int| & Size in bytes \\
\lstinline|"mtime"| & \lstinline|Int| & Last modification time (epoch milliseconds) \\
\lstinline|"mode"| & \lstinline|Int| & Unix permission bits (e.g., \texttt{0o755}) \\
\lstinline|"permissions"| & \lstinline|Int| & Same as \lstinline|"mode"| (alias) \\
\bottomrule
\end{tabular}
\end{center}

If you only need the file size, \lstinline|file_size| is a convenient
shortcut that returns an \lstinline|Int| directly:

\begin{lstlisting}[language=Lattice, caption={Quick file size check}]
let bytes = file_size("database.db")
let megabytes = bytes / (1024 * 1024)
print("Database is " + to_string(megabytes) + " MB")
\end{lstlisting}

\subsection{Copying, Renaming, and Deleting}\label{sec:copy-rename-delete}

\index{copy\_file}
\index{rename}
\index{delete\_file}

\begin{lstlisting}[language=Lattice, caption={File operations: copy, rename, delete}]
// Make a backup before modifying
copy_file("config.toml", "config.toml.bak")

// Rename (move) a file
rename("draft.md", "final.md")

// Clean up temporary files
delete_file("scratch.tmp")
\end{lstlisting}

\lstinline|copy_file| reads the source in 8\,KB chunks and writes them to the
destination, making it safe for large files. \lstinline|rename| performs an
atomic rename on the same filesystem (it delegates to the C
\texttt{rename()} function). \lstinline|delete_file| unlinks a single file;
it cannot remove directories.

\subsection{Permissions and Canonical Paths}\label{sec:chmod-realpath}

\index{chmod}
\index{realpath}

\begin{lstlisting}[language=Lattice, caption={Changing permissions and resolving paths}]
// Make a script executable
write_file("deploy.sh", "#!/bin/sh\necho 'Deploying...'\n")
chmod("deploy.sh", 0o755)

// Resolve symlinks and relative components
let canonical = realpath("./src/../README.md")
print(canonical)  // "/home/user/project/README.md"
\end{lstlisting}

\lstinline|chmod| takes a path and an integer mode. You can use octal
literals like \lstinline|0o755| for readability. \lstinline|realpath| resolves
a path to its absolute, canonical form---expanding symlinks, removing
\texttt{.} and \texttt{..} components, and verifying that the path exists.

\begin{defnbox}{Filesystem Functions at a Glance}
All filesystem functions are registered as top-level built-ins in
\filename{src/runtime.c}. They require no imports. On Emscripten (browser)
builds, all filesystem operations return errors, since there is no real
filesystem available---check the Emscripten stubs in
\filename{src/fs\_ops.c} if you are curious about the boundary.
\end{defnbox}

% ---------------------------------------------------------------------------
\section{The path Module}\label{sec:path-module}
% ---------------------------------------------------------------------------

\index{path\_join}
\index{path\_dir}
\index{path\_base}
\index{path\_ext}

File paths are strings, but they have structure. The path functions help you
manipulate that structure without resorting to brittle string concatenation.

\subsection{Joining Paths: path\_join}\label{sec:path-join}

\begin{lstlisting}[language=Lattice, caption={Joining path components}]
let dir = "/home/user/projects"
let file = "lattice/main.lat"
let full = path_join(dir, file)
print(full)  // "/home/user/projects/lattice/main.lat"
\end{lstlisting}

\lstinline|path_join| is variadic---you can pass any number of string
arguments:

\begin{lstlisting}[language=Lattice, caption={Joining multiple path segments}]
let config_path = path_join("/etc", "lattice", "config.toml")
print(config_path)  // "/etc/lattice/config.toml"
\end{lstlisting}

The implementation in \filename{src/path\_ops.c} is careful about separators.
It avoids producing double slashes (\texttt{//}) when one component ends with
a slash and the next begins with one, and it inserts a separator when neither
component has one. This means you do not have to worry about trailing slashes
on your directory paths:

\begin{lstlisting}[language=Lattice, caption={path\_join handles trailing slashes}]
// All of these produce the same result:
print(path_join("/tmp/", "file.txt"))   // "/tmp/file.txt"
print(path_join("/tmp", "file.txt"))    // "/tmp/file.txt"
print(path_join("/tmp/", "/file.txt"))  // "/tmp/file.txt"
\end{lstlisting}

\subsection{Dissecting Paths: path\_dir, path\_base, path\_ext}
\label{sec:path-dissect}

Three functions let you take a path apart:

\begin{lstlisting}[language=Lattice, caption={Dissecting a file path}]
let path = "/home/user/photos/vacation.jpg"

print(path_dir(path))   // "/home/user/photos"
print(path_base(path))  // "vacation.jpg"
print(path_ext(path))   // ".jpg"
\end{lstlisting}

\lstinline|path_dir| returns the directory portion of a path---everything
before the last slash. If the path has no slash, it returns \lstinline|"."|
(the current directory).

\lstinline|path_base| returns the filename---everything after the last
slash. If the path ends with a slash, it returns an empty string.

\lstinline|path_ext| returns the file extension, \emph{including the dot}. If
there is no extension, or if the filename starts with a dot (like
\texttt{.gitignore}), it returns an empty string.

\begin{lstlisting}[language=Lattice, caption={Edge cases in path functions}]
// No directory component
print(path_dir("readme.md"))   // "."
print(path_base("readme.md"))  // "readme.md"

// Hidden files
print(path_ext(".gitignore"))  // "" (dot-prefixed names are not extensions)

// No extension
print(path_ext("Makefile"))    // ""

// Root path
print(path_dir("/"))           // "/"
\end{lstlisting}

\subsection{Building a File Organizer}\label{sec:file-organizer}

Let's combine path functions with filesystem operations to build something
practical---a script that organizes files into directories by extension:

\begin{lstlisting}[language=Lattice, caption={Organizing files by extension}]
fn organize_by_extension(source_dir: String) {
    let files = list_dir(source_dir)

    for name in files {
        let full_path = path_join(source_dir, name)

        // Skip directories
        if is_dir(full_path) { continue }

        // Determine the target directory
        let ext = path_ext(name)
        if ext == "" { ext = "no_extension" }
        // Remove the leading dot
        let category = ext.slice(1, ext.len())

        let target_dir = path_join(source_dir, category)
        if !file_exists(target_dir) {
            mkdir(target_dir)
        }

        let dest = path_join(target_dir, name)
        rename(full_path, dest)
        print("Moved " + name + " -> " + category + "/")
    }
}

organize_by_extension("/tmp/downloads")
\end{lstlisting}

\begin{tipbox}{Use path Functions, Not String Slicing}
It is tempting to split paths on \texttt{"/"} manually, but
\lstinline|path_dir|, \lstinline|path_base|, and \lstinline|path_ext| handle
edge cases that manual string manipulation easily gets wrong: root paths,
trailing slashes, dot-prefixed filenames, and empty components. Always prefer
the path functions.
\end{tipbox}

% ---------------------------------------------------------------------------
\section{Temporary Files and Directories}\label{sec:tempfiles}
% ---------------------------------------------------------------------------

\index{tempfile}
\index{tempdir}

Sometimes you need a place to stash data that should not outlive your program.
Lattice provides two functions for this: \lstinline|tempfile| and
\lstinline|tempdir|.

\begin{lstlisting}[language=Lattice, caption={Creating temporary files and directories}]
// Create a temporary file
let tmp_path = tempfile()
print(tmp_path)  // e.g., "/tmp/lattice_a3kJ9x"

write_file(tmp_path, "intermediate results go here")
let data = read_file(tmp_path)
print(data)  // "intermediate results go here"

// Clean up when done
delete_file(tmp_path)
\end{lstlisting}

\lstinline|tempfile| creates a new file in the system's temporary directory
(typically \texttt{/tmp} on Unix or the path returned by
\texttt{GetTempPathA} on Windows) with a unique name based on the template
\texttt{lattice\_XXXXXX}. Internally, it uses \texttt{mkstemp}, which
atomically creates the file and guarantees no name collisions. The file
descriptor is immediately closed, and the path is returned as a string.

\lstinline|tempdir| works the same way but creates a directory instead:

\begin{lstlisting}[language=Lattice, caption={Using a temporary directory as a workspace}]
let workspace = tempdir()
print(workspace)  // e.g., "/tmp/lattice_Qm7kP2"

// Use the directory for intermediate files
write_file(path_join(workspace, "step1.txt"), "Phase 1 output")
write_file(path_join(workspace, "step2.txt"), "Phase 2 output")

// List what we created
let files = list_dir(workspace)
for name in files {
    print("  " + name)
}

// Clean up
for name in files {
    delete_file(path_join(workspace, name))
}
rmdir(workspace)
\end{lstlisting}

\begin{warnbox}{Temporary Files Are Not Automatically Cleaned Up}
Lattice does not use finalizers or RAII-style destructors to clean up
temporary files. If your program exits without deleting them, they will
remain on disk until the operating system's periodic cleanup runs (if ever).
Always delete your temporaries explicitly, ideally using a \lstinline|defer|
block to ensure cleanup happens even if an error occurs:
\begin{lstlisting}[language=Lattice]
let tmp = tempfile()
defer { delete_file(tmp) }
// ... work with the file ...
\end{lstlisting}
\end{warnbox}

\subsection{A Safe Processing Pipeline}\label{sec:safe-pipeline}

A robust pattern for file processing is to write output to a temporary file,
then atomically rename it to the final destination. This prevents other
programs from reading a half-written file:

\begin{lstlisting}[language=Lattice, caption={Atomic file writes with tempfile and rename}]
fn safe_write(destination: String, content: String) {
    let dir = path_dir(destination)
    let tmp = tempfile()
    defer {
        // Clean up the temp file if rename failed
        if file_exists(tmp) {
            delete_file(tmp)
        }
    }

    write_file(tmp, content)
    rename(tmp, destination)
}

safe_write("output/report.json", "{\"status\": \"complete\"}")
\end{lstlisting}

The \lstinline|defer| block ensures the temporary file is cleaned up no
matter what happens. If the \lstinline|rename| succeeds, the temporary file
no longer exists at its original path, so the \lstinline|file_exists| check
prevents a spurious error.

% ---------------------------------------------------------------------------
\section{Working with Buffers}\label{sec:buffers}
% ---------------------------------------------------------------------------

\index{Buffer}
\index{binary data}
\index{read\_file\_bytes}
\index{write\_file\_bytes}

Strings are great for text, but the world is full of binary data: images,
network packets, serialized structures, compressed archives. Lattice
represents raw byte sequences with the \lstinline|Buffer| type.

\begin{defnbox}{Buffer}
A \emph{Buffer} is a contiguous, mutable sequence of bytes (\texttt{uint8\_t}
values, ranging from 0 to 255). Buffers have a length (the number of bytes
currently stored) and a capacity (the allocated space). You can index into a
Buffer to read or write individual bytes, and use typed read/write methods to
interpret sequences of bytes as integers or floating-point values.
\end{defnbox}

\subsection{Creating Buffers}\label{sec:buffer-create}

There are several ways to create a Buffer:

\begin{lstlisting}[language=Lattice, caption={Creating Buffers}]
// Allocate a zero-filled buffer of a given size
let zeros = Buffer::new(16)
print(zeros.len())  // 16

// Create from an array of byte values
let header = Buffer::from([0x89, 0x50, 0x4E, 0x47])
print(header.len())  // 4

// Create from a string (UTF-8 bytes)
let greeting = Buffer::from_string("Hello")
print(greeting.len())  // 5
print(greeting[0])     // 72 (ASCII 'H')
\end{lstlisting}

\lstinline|Buffer::new(size)| allocates a buffer of the given size, filled
with zeros. \lstinline|Buffer::from(array)| creates a buffer from an array
of integers, truncating each value to the range 0--255.
\lstinline|Buffer::from_string(str)| copies the UTF-8 bytes of a string into
a new buffer.

\subsection{Reading and Writing Binary Files}\label{sec:binary-files}

To read a file as raw bytes, use \lstinline|read_file_bytes|:

\begin{lstlisting}[language=Lattice, caption={Reading a binary file}]
let data = read_file_bytes("image.png")
print("File size: " + to_string(data.len()) + " bytes")

// Check the PNG magic number
if data[0] == 0x89 && data[1] == 0x50 && data[2] == 0x4E && data[3] == 0x47 {
    print("Valid PNG header!")
}
\end{lstlisting}

To write a buffer to disk, use \lstinline|write_file_bytes|:

\begin{lstlisting}[language=Lattice, caption={Writing binary data to a file}]
// Create a minimal BMP file header (simplified)
flux bmp = Buffer::new(0)
bmp.push(0x42)  // 'B'
bmp.push(0x4D)  // 'M'

// Write to disk
write_file_bytes("output.bmp", bmp)
\end{lstlisting}

\subsection{Indexing and Mutation}\label{sec:buffer-indexing}

\index{Buffer!indexing}

Buffers support integer indexing for both reading and writing:

\begin{lstlisting}[language=Lattice, caption={Buffer indexing}]
flux buf = Buffer::new(4)
buf[0] = 0xFF
buf[1] = 0xCA
buf[2] = 0xFE
buf[3] = 0x00

print(buf[0])  // 255
print(buf[2])  // 254

// Values are masked to a single byte
buf[0] = 256
print(buf[0])  // 0 (only the low 8 bits are kept)
\end{lstlisting}

When you assign a value to a buffer index, Lattice masks it with
\texttt{0xFF}, keeping only the lowest 8 bits. Reading an index returns an
\lstinline|Int| in the range 0--255. Out-of-bounds access produces a runtime
error.

\subsection{The push Method and Dynamic Growth}\label{sec:buffer-push}

\index{Buffer!push}

Buffers grow dynamically when you use \lstinline|push|:

\begin{lstlisting}[language=Lattice, caption={Building a buffer byte by byte}]
flux packet = Buffer::new(0)
packet.push(0x01)      // version
packet.push(0x00)      // flags
packet.push_u16(1024)  // payload length as 16-bit value
packet.push_u32(42)    // sequence number as 32-bit value

print(packet.len())       // 8
print(packet.capacity())  // at least 8
\end{lstlisting}

\lstinline|push| appends a single byte. \lstinline|push_u16| and
\lstinline|push_u32| append multi-byte values in little-endian byte order.

\subsection{Typed Read and Write Methods}\label{sec:typed-buffer-access}

\index{Buffer!read\_u8}
\index{Buffer!write\_u8}
\index{Buffer!read\_u16}
\index{Buffer!read\_u32}
\index{Buffer!read\_u64}
\index{Buffer!read\_i8}
\index{Buffer!read\_i16}
\index{Buffer!read\_i32}
\index{Buffer!read\_i64}
\index{Buffer!read\_f32}
\index{Buffer!read\_f64}

When working with structured binary data, you need to read and write
multi-byte values at specific offsets. Buffers provide a full suite of typed
accessors:

\begin{lstlisting}[language=Lattice, caption={Typed buffer reads and writes}]
flux buf = Buffer::new(20)

// Write values at specific offsets
buf.write_u8(0, 0xFF)        // 1 byte at offset 0
buf.write_u16(1, 0xCAFE)     // 2 bytes at offset 1
buf.write_u32(3, 0xDEADBEEF) // 4 bytes at offset 3

// Read them back
print(buf.read_u8(0))   // 255
print(buf.read_u16(1))  // 51966  (0xCAFE)
print(buf.read_u32(3))  // 3735928559  (0xDEADBEEF)
\end{lstlisting}

The full set of typed methods is:

\begin{center}
\begin{tabular}{llp{7cm}}
\toprule
\textbf{Read} & \textbf{Write} & \textbf{Description} \\
\midrule
\lstinline|read_u8(offset)| & \lstinline|write_u8(offset, val)| & Unsigned 8-bit integer (1 byte) \\
\lstinline|read_u16(offset)| & \lstinline|write_u16(offset, val)| & Unsigned 16-bit integer (2 bytes) \\
\lstinline|read_u32(offset)| & \lstinline|write_u32(offset, val)| & Unsigned 32-bit integer (4 bytes) \\
\lstinline|read_u64(offset)| & \lstinline|write_u64(offset, val)| & Unsigned 64-bit integer (8 bytes) \\
\lstinline|read_i8(offset)| & --- & Signed 8-bit integer \\
\lstinline|read_i16(offset)| & --- & Signed 16-bit integer \\
\lstinline|read_i32(offset)| & --- & Signed 32-bit integer \\
\lstinline|read_i64(offset)| & \lstinline|write_i64(offset, val)| & Signed 64-bit integer \\
\lstinline|read_f32(offset)| & --- & IEEE 754 single-precision float (4 bytes) \\
\lstinline|read_f64(offset)| & --- & IEEE 754 double-precision float (8 bytes) \\
\bottomrule
\end{tabular}
\end{center}

All multi-byte reads and writes use \emph{little-endian} byte order, matching
the native byte order on x86 and ARM systems.

\begin{warnbox}{Bounds Checking}
All typed read and write methods perform bounds checking. Reading a
\lstinline|u32| at offset 5 from a 7-byte buffer will produce a runtime error,
because the read would extend past the end of the buffer (offset 5 + 4 bytes
= offset 9, which exceeds length 7). Always verify your offsets when parsing
binary formats.
\end{warnbox}

\subsection{Slicing, Filling, and Converting}\label{sec:buffer-utility}

\index{Buffer!slice}
\index{Buffer!fill}
\index{Buffer!clear}
\index{Buffer!resize}
\index{Buffer!to\_string}
\index{Buffer!to\_array}
\index{Buffer!to\_hex}

Buffers provide several utility methods for common operations:

\begin{lstlisting}[language=Lattice, caption={Buffer utility methods}]
let data = Buffer::from_string("Hello, Lattice!")

// Slice: extract a sub-buffer (start, end)
let sub = data.slice(0, 5)
print(sub.to_string())  // "Hello"

// to_hex: get a hex string representation
let header = Buffer::from([0xDE, 0xAD, 0xBE, 0xEF])
print(header.to_hex())  // "deadbeef"

// to_array: convert to an array of integers
let bytes = header.to_array()
print(bytes)  // [222, 173, 190, 239]

// fill: set all bytes to a value
flux scratch = Buffer::new(8)
scratch.fill(0xFF)
print(scratch[3])  // 255

// clear: reset length to zero
scratch.clear()
print(scratch.len())  // 0

// resize: change the buffer length
flux resizable = Buffer::new(4)
resizable.resize(8)
print(resizable.len())  // 8
\end{lstlisting}

\lstinline|slice(start, end)| returns a \emph{new} buffer containing the
bytes from index \texttt{start} up to (but not including) \texttt{end}.
\lstinline|to_string| interprets the buffer bytes as UTF-8 and returns a
\lstinline|String|. \lstinline|to_hex| produces a lowercase hexadecimal
representation. \lstinline|to_array| returns an \lstinline|Array| of
\lstinline|Int| values, one per byte.

\subsection{Parsing a Binary File Format}\label{sec:binary-format-example}

Let's put everything together by parsing a WAV audio file header. The WAV
format starts with a well-defined 44-byte header:

\begin{lstlisting}[language=Lattice, caption={Parsing a WAV file header}]
fn parse_wav_header(path: String) {
    let data = read_file_bytes(path)
    if data.len() < 44 {
        print("File too small to be a valid WAV")
        return
    }

    // Bytes 0-3: "RIFF" marker
    let riff = data.slice(0, 4).to_string()
    if riff != "RIFF" {
        print("Not a RIFF file")
        return
    }

    // Bytes 4-7: file size minus 8
    let file_size = data.read_u32(4) + 8

    // Bytes 8-11: "WAVE" format
    let format = data.slice(8, 12).to_string()
    if format != "WAVE" {
        print("Not a WAVE file")
        return
    }

    // Bytes 22-23: number of channels
    let channels = data.read_u16(22)

    // Bytes 24-27: sample rate
    let sample_rate = data.read_u32(24)

    // Bytes 34-35: bits per sample
    let bits_per_sample = data.read_u16(34)

    print("WAV File: " + path)
    print("  Channels:       " + to_string(channels))
    print("  Sample rate:    " + to_string(sample_rate) + " Hz")
    print("  Bits per sample: " + to_string(bits_per_sample))
    print("  File size:      " + to_string(file_size) + " bytes")
}

parse_wav_header("recording.wav")
\end{lstlisting}

This example demonstrates the core workflow for binary file parsing: read the
file as bytes, check magic numbers, and use typed reads at known offsets to
extract structured fields.

\begin{notebox}{Buffers Are Value Types}
Like structs in Lattice, Buffers are \emph{value types}---assigning a buffer
to a new variable or passing it to a function creates a copy. If you
modify the copy, the original is unaffected. This eliminates a whole category
of aliasing bugs but means you should be mindful of copying large buffers
unnecessarily. For more on value semantics, see \cref{ch:values-types}.
\end{notebox}

% ---------------------------------------------------------------------------
\section*{Exercises}\label{sec:filesystem-exercises}
% ---------------------------------------------------------------------------

\begin{enumerate}
    \item \textbf{Disk Usage Counter.} Write a function
    \lstinline|disk_usage(dir: String) -> Int| that recursively calculates
    the total size in bytes of all files in a directory tree. Use
    \lstinline|list_dir|, \lstinline|is_dir|, \lstinline|is_file|, and
    \lstinline|file_size|.

    \item \textbf{Find and Replace.} Write a program that takes a directory
    path, a search string, and a replacement string as arguments. It should
    find all \texttt{.txt} files in the directory using \lstinline|glob|,
    read each one, replace all occurrences of the search string, and write
    the result back. Use the safe-write pattern from
    \cref{sec:safe-pipeline} to avoid corrupting files.

    \item \textbf{Hex Dump.} Write a function
    \lstinline|hex_dump(path: String, count: Int)| that reads the first
    \lstinline|count| bytes of a file using \lstinline|read_file_bytes| and
    prints them in a classic hex dump format: 16 bytes per line, showing
    both hexadecimal values and printable ASCII characters.

    \item \textbf{File Deduplication.} Write a program that scans a
    directory, reads every file into a Buffer, computes a checksum (you can
    use the \lstinline|to_hex| method on the bytes for a simple hash, or
    wait until \cref{ch:time-crypto-os} for \lstinline|sha256|), and
    reports any duplicate files.

    \item \textbf{Binary Concatenator.} Write a function that takes an array
    of file paths and a destination path, reads each file as bytes, and
    concatenates all the buffers into a single file. Pay attention to
    building the output buffer efficiently.
\end{enumerate}

% ---------------------------------------------------------------------------
\section*{What's Next}\label{sec:filesystem-next}
% ---------------------------------------------------------------------------

Now that we can read and write files, the natural next question is: what
\emph{format} should those files be in? In the next chapter, we will explore
Lattice's built-in support for four of the most common data formats---JSON,
TOML, YAML, and CSV---and see how to parse, generate, and round-trip
structured data between them.
